\documentclass[a4paper,10pt]{article}
\usepackage[utf8]{inputenc}

% ----  Useful packages % ---- 
\usepackage{amsmath}
\usepackage{graphicx}
\usepackage{amsfonts}
\usepackage{amsthm}
\usepackage{amssymb}
\usepackage{makecell}
\usepackage{tikz}
\usepackage{pgfplots}
% ----  Useful packages % ---- 

\usepackage{wrapfig}
\usepackage{caption}
\usepackage{subcaption}
\usepackage{hyperref}
\hypersetup{
    colorlinks,
    citecolor=black,
    filecolor=black,
    linkcolor=black,
    urlcolor=black
}

\graphicspath{ {./images/} }

% ---- Set page size and margins replace ------
\usepackage[letterpaper,top=2cm,bottom=2cm,left=3cm,right=3cm,marginparwidth=1.75cm]{geometry}
% ---- Set page size and margins replace ------

% ------- NOTA ------
\theoremstyle{remark}
\newtheorem{note}{Note}[subsection]
% ------- NOTA ------

% ------- OSSERVAZIONE ------
\theoremstyle{definition}
\newtheorem{observation}{Osservazione}[subsection]
% ------- OSSERVAZIONE ------

% ------- DEFINIZIONE ------
\theoremstyle{plain}
\newtheorem{definition}{Definizione}[subsection]
% ------- DEFINIZIONE ------

% ------- ESEMPIO ------
\theoremstyle{definition}
\newtheorem{example}{Esempio}[subsection]
% ------- ESEMPIO ------

% ------- DIMOSTRAZIONE ------
\theoremstyle{definition}
\newtheorem{demostration}{Dimotrazione}[subsection]
% ------- DIMOSTRAZIONE ------

% ------- TEOREMA ------
\theoremstyle{definition}
\newtheorem{theorem}{Teorema}[subsection]
% ------- TEOREMA ------

% ------- COROLLARIO ------
\theoremstyle{plain}
\newtheorem{corollaries}{Corollario}[theorem]
% ------- COROLLARIO ------

% ------- PROPOSIZIONE ------
\theoremstyle{plain}
\newtheorem{proposition}{Proposizione}[subsection]
% ------- PROPOSIZIONE ------

% ---- Footer and header ---- 
\usepackage{fancyhdr}
\pagestyle{fancy}
\fancyhf{}
\fancyhead[LE,RO]{A.A 2023-2024}
\fancyhead[RE,LO]{Calcolo Numerico}
\fancyfoot[RE,LO]{\rightmark}
\fancyfoot[LE,RO]{\thepage}

\renewcommand{\headrulewidth}{.5pt}
\renewcommand{\footrulewidth}{.5pt}
% ---- Footer and header ---- 

% ----  Language setting ---- 
\usepackage[italian, english]{babel}
% ----  Language setting ---- 

\usepackage{listings}
\usepackage{color}

\definecolor{dkgreen}{rgb}{0,0.6,0}
\definecolor{gray}{rgb}{0.5,0.5,0.5}
\definecolor{mauve}{rgb}{0.58,0,0.82}

\lstset{frame=tb,
  language=C,
  aboveskip=3mm,
  belowskip=3mm,
  showstringspaces=false,
  columns=flexible,
  basicstyle={\small\ttfamily},
  numbers=none,
  numberstyle=\tiny\color{gray},
  keywordstyle=\color{blue},
  commentstyle=\color{dkgreen},
  stringstyle=\color{mauve},
  breaklines=true,
  breakatwhitespace=true,
  tabsize=3
}

\title{\textbf{Calcolo Numerico}}
\author{Realizzato da: Giuntoni Matteo e Filippo Ghirardini}
\date{A.A. 2023-2024}

\begin{document}
\begin{titlepage} %crea l'enviroment
	\begin{figure}[t] %inserisce le figure
		\centering\includegraphics[width=0.98\textwidth]{marchio_unipi_pant541.png}
	\end{figure}
	\vspace{20mm}
	
	\begin{Large}
		\begin{center}
			\textbf{Dipartimento di Informatica\\ Corso di Laurea Triennale in Informatica\\}
			\vspace{20mm}
			{\LARGE{Corso a Libera Scelta - 6 CFU}}\\
			\vspace{10mm}
			{\huge{\bf Computer Graphics}}\\
		\end{center}
	\end{Large}
	
	
	\vspace{36mm}
	%minipage divide la pagina in due sezioni settabili
	\begin{minipage}[t]{0.47\textwidth}
		{\large{\bf Professore:}\\ \large{Prof. }}
	\end{minipage}
	\hfill
	\begin{minipage}[t]{0.47\textwidth}\raggedleft
		{\large{\bf Autore:}\\ \large{Filippo Ghirardini}}
	\end{minipage}
	
	\vspace{25mm}
	
	\hrulefill
	
	\vspace{5mm}
	
	\centering{\large{\bf Anno Accademico 2023/2024 }}
	
\end{titlepage}
\tableofcontents
\newpage
\maketitle
\begin{center}
    \vspace{-20pt}
    \rule{11cm}{.1pt} 
\end{center}
\newpage
\section{Aritmetica di Macchina}
Per una macchina la scrittura $(x + y) + z \neq x + (y + z)$. Vediamo dunque che ci sono alcuni punti focali da considerare
per far si che una macchina funzioni correttamente:

\begin{itemize}
    \item Trovare uno standard per come memorizzare i numeri.
    \item Trovare uno standard per come manipolare i numeri.
\end{itemize}

\hspace{-15pt}Da questi due punti possiamo ricondurci ad un solo problema, come andare a rappresentare i numeri.

\subsection{Teorema di rapresentazione}
\begin{theorem}
    Dato $x \in \mathbb{R}, x \neq 0$ esistono e sono univocamente determinati.
    \begin{enumerate}
        \item un intero $p \in \mathbb{Z}$ detto esponente della rappresentazione.
        \item una successione di numeri naturiali $\{d_i\}_{i\geq 1}$ con $d_i \neq 0, 0 \leq d_i \leq B - 1$ e $d_i$
        non definitvamente uguali a \(B - 1\), dette cire della rappresentazione; tali per cui si ha
    \end{enumerate}
    \begin{equation}
        x = sign(x)B^p \sum_{i=1}^{+\infty}d_i B^{-1}.
    \end{equation}
\end{theorem}

\hspace{-15pt}Andiamo ora ad analizzare il significato di questo teorema. Esso descrive quella che viene chiamata rappresentazione
in virgola mobile, in quanto l'esponente \(p\) on è determinato in modo da avere la parte intera nulla. Le cose da considerare
in questo teorema sono:
\begin{itemize}
    \item La condizione \(d_i \neq 0\) e \(d_i\) non definitivamente uguale a \(B - 1\) sono introdotte per garantire
    l'unicità delle rappresentazioni. Ad esempio:
    \[B = 10 \text{ abbiamo } 1 = +10^1 (1 \cdot 10^{-1}) = +10^2 (0 \cdot 10^{-1} + 1 \cdot 10^{-1})\]
    Quindi due rappresentazioni diverse per lo stesso numero, però considerando le condizioni scritte sopra la seconda non risulta
    accettabile perché la prima cifra è nulla.
    \item Il caso \(x=0\) non ammette rapresentazione normalizzata. Questa casistica viene trattata dalla macchina in un modo
    particolare, per questo abbiamo la condizione \(x\neq \).
    \item Questa rapresentazione si estende anche all'insieme dei numeri complessi del tipo \(z = a + ib\), utilizzando una rapresentazione
    come coppie di numeri reali del tipo \((a,b)\).
\end{itemize}

Possiamo dedurre che visto che stiamo lavornano con registri di meoria di un calcolatore con memoria a numero finito, anche la quantità
di cifre rapresentabili saranno a numero finito esso vinene chiamato \textbf{insieme dei numeri di macchina}.\\

Dal teorema di rapresentazione in base di un numero reale può avvenire assegnando delle posizioni di meoria per il segno, 
per l'esponente e per le cifre della rappresentazione.

\begin{definition}[Insieme die numeri di macchina]
    Si definisce l'insieme dei nuermi di macchina in rappresentazione floating point con t cifre, base B e range \(-m, M\) l'insieme dei numeri reali.
    \[\mathbb{F}(B, t, m, M) = \{0\} \cup \{s \in \mathbb{R} \::\: x = sign(x) B^p \sum_{i=1}^{t} d_i B^{-1}, 0\leq d_i \leq B - 1, d_1 \neq 0, -m \leq p \leq M\}\]
\end{definition}

\hspace{-15pt}Si osserva in questa definizione che:
\begin{itemize}
    \item L'insieme \(\mathbb{F}\) ha cardinalità finita \(N = 2B^{t-1}(B-1)(M + m + 1) + 1\).
    \item L'insieme dei numeri di macchina \(\mathbb{F}(B, t, m, M)\) è simmetrico rispetto all'origine.
    \item Possiamo definire \(\Omega = B^M (B-1) \sum_{i=1}^{t} b^{-1}\) come il più grande numero macchina e \(\omega = B^{-m} B^{-1}\) come
    invece il più piccolo.
    \item Posto un \(x = B^P \sum_{i=1}^{t}d_i B^{-1}\) possiamo definire il suo successiovo numero di macchina come
    \(y = B^p(\sum_{i=1}^{t-1}d_i B^{-1} + (d_t + 1)B^{-t})\). Da qui vediamo che la distanza \(y - x = B^p - t\) porta
    i numeri ad essere non equispaziali fra di loro, quindi la distanza aumento con l'avicinarci a \(\Omega\)
\end{itemize}

\begin{example}
    Facciamo ora un esempio in cui andiamo a rappresentare il numero successivo di \(x = B^p \sum_{t}^{i=1} d_i B^{-1}\). Esso
    si può scrivere come \(y = B^p \bigg( \sum_{i=1}^{t-1} d_i B^{-1} + (d_t + 1)B^{-t} \bigg)\).\\
    Mentre si può scrivere la distanza fra questi due valori come \(y - x = B^p - t\).
\end{example}

\hspace{-15pt}E' stato fissato uno standard IEEE 754 fra gli anni 70/80, questo standard dice che, visto ci sono macchine che hanno metodi di rappresentazione
diversi bisogna fissare un standard, esso appunto dice che \(B = 2\) ed i registri sono a 32 o 64 bit.\\\\
Questa rapresentazione ha uno svantaggio che può sembrare minimo ma non lo è, lo 0 si rappresenta due volte con \(-0, +0\). Per ovviare
a questo problema si è andato ad abbandonare questa rappresentazione in esponenti ma si rapprensentato i numeri nel seguente modo:
\(p_1 2^0 + p_1 2^1 + \dots + p_11 s^10 \) che rappresentano numeri da 0 a \(2^11 - 1\) quindi 2047 numeri, mentre lo 0 si può scrivere come:
\begin{itemize}
    \item O tenendo tutti i valori a 0
    \item Oppure tendendo tutti i valori a 1
\end{itemize}

\hspace{-15pt}In entrambi i casi abbiamo un range di valori che va da \([-1022, 1024]\). A questo punto ho \(2^{P-1022}\) numeri che la 
macchina rappresneta come \(\pm 2^{P - 1022}(0.1d_1 \dots d_{52})\).\\
Impostando questo standard abbiamo \(\Omega = 2^{1024}(01 \dots 1)_2\) e \(\omega = 2^{-1022} (101)_2\).\\

\begin{observation}
    Quando \(p=0\) abbiamo i numeri che si trovano nella porzione della retta dei numeri che è compresa fra \(-\omega\) \(\omega\)
    e possiamo qui avere anche tuti 0 e quindi si introduce il caso dei numeri denormalizzati.
\end{observation}

\hspace{-15pt}Se abbiamo l'esponante uguale a tutti 1, la convensione è che tutte le cifre della mantissa sono tutti uguali a 0/1
questo numero indica il \(\pm \infty\) altrimenti sta a signifiare NaN (not a number). Questi valori ci permettono di gestire
forme indeterminate.

\subsection{Errore di rappresentazione}
Quando si va a rappresentare un numero reale non nullo \(x \in \mathbb{R}\) e con \(x \neq 0\) si può andare a commettere degli 
errori di rappresentazione detto anche \textbf{errore relativo di approssimazione}, e di definisce come, prendendo un \(\tilde{x} \in \mathbb{F}(B, t, m, M)\)
\[\epsilon_x = \frac{\tilde{x} - x}{x} = \frac{\eta x}{x}, x \neq 0\]

\hspace{-15pt}Definiamo \(|\epsilon_x| = \big | \frac{\tilde{x} - x}{x}\big | \leq \frac{B^{P-t}}{| x|} \leq \frac{B^{P - t}}{B^{P - 1}} = B^{1 - t} = u\)
la \(u\) è definita come \textbf{precisione di macchina}.\\

\hspace{-15pt}Andiamo inoltre a definire le conidizioni di underflow e overflow. Dato un \(x \in \mathbb{R}, x \neq 0\) abbiamo che:
\begin{enumerate}
    \item Se \(|x| < \omega\) o \(|x| > \Omega\) overflow. In questo caso si va ad associare il \(+\infty\).
    \item Se invece \(\omega \leq |x| \leq \Omega\) abbiamo underflow. In questo caso allora prendiamo una \(x = B^p \sum_{i = 1}^{\infty}d_i B^{-1}
    \longrightarrow B^p \sum_{i=1}^{t}d_i B^{-1} = \tilde{x}\) che è una approssimazione 
\end{enumerate}
% !TeX spellcheck = it_IT
\newpage
\section{Calcolo degli errori}
Supponiamo di avere una funzione \(f: [a,b] \to \mathbb{R}\) e \(f \neq 0\), per andare a calcolare questa funzione come già visto usiamo un algoritmo
che esprime tale valore come risutlato di una sequanza di operazioni aritmentiche. Questa rappresentazione come abbiamo già potuto
verificare con esempi produce degli errori di approssimazione. Questi errori possono essere suddivisini in 3 tipologie.

\begin{definition}[Errore inerente o inevitabile]
    Si dice errore inerente o inevitabile generato nel calcolo di \(f(x) \neq 0\) la quantità:
    \[\epsilon_{in} = \frac{f(\tilde{x}) - f(x)}{f(x)}\]
\end{definition}

\begin{definition}[Errore algoritmico]
    Si dice errore algoritmico generato nel calcolo di \(f(x) \neq 0\) la quantità:
    \[\epsilon_{alg} = \frac{g(\tilde{x}) - f(\tilde{x})}{f(\tilde{x})}\]
\end{definition}

\begin{definition}[Errore totale]
    Si dice errore algoritmico totale nel calcolo di \(f(x) \neq 0\) mediante l'algoritmo specificato da g la quantità:
    \[\epsilon_{tot} = \frac{g(\tilde{x}) - f(x)}{f(x)}\]
\end{definition}

\begin{observation}
    Vediamo che se \(|\epsilon_{in}|\) è grande il problema si definisce \textbf{problema mal condizionato}. Mentre se
    \(|\epsilon_{alg}\) è grande l'algoritmo si dice che \textbf{algoritmo è numericamente instabile}.
\end{observation}
% !TeX spellcheck = it_IT
\newpage
\section{Matrici}
%TODO recuper

\subsection{Norme vettoriali}
Sono uno strumento che ci permette di definire una distanza tra due vettori.
\begin{definition}[Norma vettoriale]
	È una funzione del tipo $f:\mathbb{R}^n \to \mathbb{R}$ che deve soddisfare tre proprietà:
	\begin{enumerate}
		\item $f(x)\geq 0 \wedge f(x)=0 \Leftrightarrow x=0$
		\item $f(\alpha x)=\lvert \alpha \rvert f(x) \quad \forall \alpha \in \mathbb{R} \quad \forall x \in \mathbb{R}^n$
		\item \textbf{Disuguaglianza triangolare}: $f(x+y)\leq f(x)+f(y) \quad \forall x,y \in \mathbb{R}^n$
	\end{enumerate}
	e la indichiamo come
	\begin{equation*}
		f(x) = \lVert x \rVert
	\end{equation*}
\end{definition}

Detto questo possiamo definire una distanza come:
\begin{equation}
	dist(x,y)=\lVert  x-y \rVert
\end{equation}
Le tre proprietà ci danno alcune informazioni sulla distanza:
\begin{enumerate}
	\item La distanza deve essere non negativa e valere $0$ solo se i due vettori coincidono 
	\item La distanza tra $x$ e $y$ deve essere uguale a quella tra $y$ e $x$
	\item $\lVert x-y \rVert = \lVert(x-a)+(a-y)\rVert \leq \lVert x-a \rVert + \lVert a-y \rVert$
\end{enumerate}

\begin{definition}[Distanza euclidea - Norma 2]
	\begin{equation}
		f(x) = \sqrt{x_1 ^ 2 + x_2 ^ 2 + \ldots + x_n ^ 2} > \lVert  x \rVert_2
	\end{equation}
\end{definition}
\begin{definition}[Norma infinito]
	\begin{equation}
		f(x)=max\lvert x_1 \rvert = \lVert x \rVert_{\infty}
	\end{equation}
\end{definition}
\begin{definition}[Norma 1]
	\begin{equation}
		f(x)=\sum_{i=1}^{n}\lvert x_i \rvert = \lVert x\rVert_1
	\end{equation}
\end{definition}

\begin{example}
	Prendiamo due vettori
	\begin{equation*}
		\begin{Bmatrix}
			1 \\
			1
		\end{Bmatrix} \quad
		\begin{Bmatrix}
			2 \\
			-2
		\end{Bmatrix}
	\end{equation*}
	e calcoliamo le varie norme:
	\begin{equation*}
		\lVert x \rVert_2 = \sqrt{2} \quad \lVert x \rVert_\infty = 1 \quad \lVert x \rVert_1 = 2
	\end{equation*}
	%TODO Finisci
\end{example}

\begin{definition}[Norma matriciale]
		È una funzione del tipo $f:\mathbb{R}^{n\times n} \to \mathbb{R}$ che deve soddisfare tre proprietà:
	\begin{enumerate}
		\item $f(A)\geq 0 \wedge f(A)=0 \Leftrightarrow A=0$
		\item $f(\alpha A)=\lvert \alpha \rvert f(x) \quad \forall \alpha \in \mathbb{R} \quad \forall A \in \mathbb{R}^{n\times n}$
		\item \textbf{Disuguaglianza triangolare}: $f(A+B)\leq f(A)+f(B) \quad \forall A,B \in \mathbb{R}^{n\times n}$
		\item $f(A\cdot B) \leq f(A) f(B)$
	\end{enumerate}
	e la indichiamo come
	\begin{equation*}
		f(A) = \lVert A \rVert
	\end{equation*}
\end{definition}
% !TeX spellcheck = it_IT
\newpage
\section{Condizionamento}
Studiare il condizionamento di un problema in forma $Ax=b$ significa chiedersi di quanto cambia la soluzione perturbando di poco $A$ e $B$.
\subsection{Metodi diretti}
Dato un sistema lineare $Ax=b$ le soluzioni possono essere trovate tramite
\begin{equation*}
	x_i = f_i(a_{11}, \ldots, a_{1n}, b1, \ldots, b_n) \quad i:1\ldots n
\end{equation*}
Si può partire studiando la forma di $A$ per cercare dei sistemi risolvibili facilmente.
\subsubsection{Matrice diagonale}
Il primo caso è quando $A$ è una matrice \textbf{diagonale}:
\begin{equation*}
	a_{ij}=0 \Longleftarrow i \neq j 
\end{equation*}
Il determinante di una matrice diagonale è $det(A)=\prod_{i=1}^{n} a_i$ e il determinante è diverso da $0$ solo se tutti gli elementi della diagonale lo sono.
Possiamo riscrivere la matrice in un sistema di equazioni lineari:
\begin{equation*}
	Ax=b \Leftrightarrow \left\{ \begin{array}{lr}
		a_1x_1 = b_1 \Leftrightarrow x1 = \frac{b1}{a1}\\
		\vdots\\
		a_nx_n = b_n \Leftrightarrow x_n=\frac{b_n}{a_n}
	\end{array} \right.
\end{equation*}
Questo sistema lo risolviamo in $O(n)$. Dato però che ridurre una matrice normale in una diagonale è un processo complesso, non conviene farlo.
\subsubsection{Matrice triangolare}
Una matrice è \textbf{triangolare inferiore} quando $a_{ij}=0$ per $j>i$, mentre è \textbf{triangolare superiore} quando $a_{ij}=0$ per $i>j$.\\
Per calcolare il determinante di una matrice triangolare superiore con Laplace facciamo:
\begin{equation*}
	A=\begin{bmatrix}
		a_{11} & 0 & \cdots & 0 \\
		\vdots & \ddots & 0 & \vdots \\
		\vdots & \ddots& \ddots & 0\\
		a_{n1} & \cdots &\cdots & a_{nn}
	\end{bmatrix}
\end{equation*}
Per risolvere $Ax=b$ con una matrice triangolare, vediamo che il sistema associato ci porta ad avere prima un equazione con tutte le incognite fino all'ultima con una sola incognita. È intuitivo che per calcolarne la soluzione è sufficiente partire dall'ultima ed eseguire una \textbf{sostituzione all'indietro}.
\begin{lstlisting}[language=Matlab]
	function[x]=backward_substitution(a,b)
		n=length(b);
		x = zeros(n,1);
		for k=n:-1:1
			s=0;
			for j=k+1:n
				s=s+a(k,j)*x(j);
			end
			x(k)=(b(k)-s)/a(k,k);
		end
\end{lstlisting}
Dato che l'operazione moltiplicativa in macchina richiede leggermente più lunga di quella della somma, è sufficiente considerare le prime per calcolare la complessità di questo programma.\\ La complessità al caso pessimo è quindi $\frac{n \cdot (n+1)}{2} = O(n^2) = \frac{n^2}{2}+O(n)$.
\begin{example}[Matrice bidiagonale superiore]
	Supponiamo di avere la seguente matrice bidiagonale superiore:
	\begin{equation*}
		A=\begin{bmatrix}
			a_{11} & a_{12} & 0 & 0\\
			0 & \ddots & \ddots & 0 \\
			0 & 0 & \ddots & a_{n-1 n} \\
			0 & 0 & 0 & a_{nn}
		\end{bmatrix}  = 
		\begin{bmatrix}
			a_1 & b_1 & 0 & 0\\
			0 & \ddots & \ddots & 0\\
			0 & 0 & \ddots & b_{n-1}\\
			0 & 0 & 0 & a_n
		\end{bmatrix}
	\end{equation*}
	In questo caso il codice si può cambiare sostituendo il secondo ciclo con l'operazione
	\begin{lstlisting}[language=Matlab]
		s=a(k,k+1)*x(k+1);
	\end{lstlisting}
	e la complessità diventa $O(n)$.
\end{example}
% !TeX spellcheck = it_IT
\section{Metodi iterativi}
Una classe di metodi che permette di risolvere sistemi lineari è quella dei metodi iterativi. L'idea è di costruire una successione di vettori partendo dalla matrice che convergano alla soluzione del sistema lineare.
\begin{equation*}
	\{X_k\}_{k\in \mathbb{N}}
\end{equation*}
\begin{observation}
	Non potendo generare infiniti termini, ad un certo punto ci sarà bisogno di fermarsi quando si pensa di essere sufficientemente vicini alla soluzione.
\end{observation}
Diciamo che
\begin{equation}
	\lim_{k\to +\infty}x_k \Leftrightarrow \lim_{k\to +\infty} \lvert\lvert x_k - x\rvert\rvert_{\infty} = 0
\end{equation}
\noindent Questo è vero perché:
\begin{equation*}
	\forall j = 1, \ldots, n \quad 0 \leq \lvert x^{(k)}_j - x_j \rvert \leq \lvert\lvert x^{(k)} - x\rvert\rvert_{\infty}
\end{equation*}

Abbiamo una matrice invertibile $A$. Supponiamo di scomporre la matrice come:
\begin{equation}
	A=M-N
\end{equation}
con l'assunzione che anche $M$ sia invertibile ($det(M)=0$). A questo punto possiamo dire che:
\begin{equation*}
	Ax=b \Leftrightarrow (M-N)x = b\Leftrightarrow Mx=Nx+b \Leftrightarrow x=M^{-1}Nx + M{-1}b \Leftrightarrow x=Px+q
\end{equation*}
Dovendo trovare la soluzione di $x=Px+q$, possiamo scegliere un vettore iniziale $x_0 \in \mathbb{R}^n$ e costruirci la successione di vettori
\begin{equation}
	\label{equation:succ}
	x_{k+1}=Px_k + q
\end{equation}
Se questa successione converge ho trovato la soluzione del sistema lineare.
\begin{observation}
	Nell'implementazione pratica non userò mai l'equazione \ref{equation:succ} ma invece
	\begin{equation}
		Mx_{k+1}=Nx_k + b
	\end{equation}
\end{observation}

\begin{example}
	Prendiamo
	\begin{equation*}
		\begin{bmatrix}
			2 & 1 \\
			1 & 2
		\end{bmatrix} \quad
		b = \begin{bmatrix}
			0 \\ 0
		\end{bmatrix} \quad
		Ax=b \Leftrightarrow x=\begin{bmatrix}
			0 \\ 0
		\end{bmatrix}
	\end{equation*}
	Partiamo definendo le due matrici per la scomposizione:
	\begin{equation*}
		M=\begin{bmatrix}
			2 & 0 \\ 0 & 2
		\end{bmatrix} \quad
		N = \begin{bmatrix}
			0 & -1 \\ -1 & 0
		\end{bmatrix}
	\end{equation*}
	e calcolando la matrice $P$
	\begin{equation*}
		P = M^{-1}N = \begin{bmatrix}
			\frac{1}{2} & 0 \\ 0 & \frac{1}{2}
		\end{bmatrix}\begin{bmatrix}
		0 & -1 \\ -1 & 0
		\end{bmatrix} = \begin{bmatrix}
		0 & -\frac{1}{2} \\ -\frac{1}{2} & 0
		\end{bmatrix}
	\end{equation*}
	e il vettore $q$
	%TODO
	Iniziamo il metodo iterativo:
	\begin{equation*}
		\begin{split}
			& x^{k+1} = Px^{(k)} = \begin{bmatrix}
				0 & -\frac{1}{2} \\ -\frac{1}{2} & 0
			\end{bmatrix} x^{(k)} \\
			& x^{(1)} =  \begin{bmatrix}
				0 & -\frac{1}{2} \\ -\frac{1}{2} & 0
			\end{bmatrix} \begin{bmatrix}
			x \\ y
			\end{bmatrix} = \begin{bmatrix}
			-\frac{1}{2} y \\ -\frac{1}{2} x
			\end{bmatrix} \\
			& x^{(2)} =  \begin{bmatrix}
				0 & -\frac{1}{2} \\ -\frac{1}{2} & 0
			\end{bmatrix} \begin{bmatrix}
				-\frac{1}{2}y \\ -\frac{1}{2}x
			\end{bmatrix} = \begin{bmatrix}
				\frac{1}{4} x \\ \frac{1}{4} y
			\end{bmatrix}
		\end{split}
	\end{equation*}
	e notiamo che la successione tende a 
	\begin{equation*}
		\begin{bmatrix}
			0 \\ 0
		\end{bmatrix}
	\end{equation*}
	Se prendiamo invece
	\begin{equation*}
		M=\begin{bmatrix}
			0 & 1 \\ 1 & 0
		\end{bmatrix} \quad
		N = \begin{bmatrix}
			-2 & 0 \\ 0 & -2
		\end{bmatrix}
	\end{equation*} con
	\begin{equation*}
		P = M^{-1}N = \begin{bmatrix}
			0 & 1 \\ 1 &0
		\end{bmatrix}\begin{bmatrix}
			-2 & 0 \\ 0 & -2
		\end{bmatrix} = \begin{bmatrix}
			0 &-2 & -2 & 0
		\end{bmatrix}
	\end{equation*}
	abbiamo che la successione \textbf{diverge}.
\end{example}
\begin{definition}
	Un metodo iterativo è convergente se
	\begin{equation}
		\forall x^{(0)} \in \mathbb{R}^n \quad x_k \to x
	\end{equation}
	Ovvero se per ogni vettore di partenza scelto, il metodo converge.
\end{definition}

\begin{theorem}
	Dato $x^{(k+1)} = Px^{(k)}+q$, se
	\begin{equation}
		\exists \lvert\lvert\cdot\rvert\rvert \text{ tc } \lvert\lvert P \rvert\rvert <1
	\end{equation}
	allora il metodo è convergente.
\end{theorem}

\begin{example}
	Data una matrice
	\begin{equation*}
		A = \begin{bmatrix}
			3 & -1 & \\ 
			-1 & \ddots & \ddots \\
			 & \ddots & \ddots & -1\\
			 & & -1 & 3
		\end{bmatrix}
	\end{equation*}
	definiamo le matrici per la scomposizione
	\begin{equation*}
		M = \begin{bmatrix}
			3 \\
			& \ddots \\
			& & 3
		\end{bmatrix} \quad
		N = \begin{bmatrix}
			0 & 1 \\
			1 & \ddots & \ddots \\
			& \ddots & \ddots & 1 \\
			& & 1 &0
		\end{bmatrix}
	\end{equation*}
	Per capire se il metodo è convergente dobbiamo calcolare la norma infinito di $P$:
	\begin{equation*}
		M^{-1} N = \begin{bmatrix}
			0 & \frac{1}{3} \\
			\frac{1}{3} & \ddots & \ddots \\
			& \ddots & \ddots & \frac{1}{3}\\
			& & \frac{1}{3} & 0
		\end{bmatrix} \quad
		\lvert\lvert P \rvert\rvert_{\infty}= \frac{1}{3} + \frac{1}{3} = \frac{2}{3} < 1
	\end{equation*}
	Il problema di questo metodo iterativo è che ha complessità $O(n)$ per ogni iterazione e non è quindi competitivo con l'eliminazione Gaussiana.
\end{example}

\begin{example}
	Data una matrice
	\begin{equation*}
		A = \begin{bmatrix}
			T & I & \\ 
			-I & \ddots & \ddots \\
			& \ddots & \ddots & I\\
			& & -I & T
		\end{bmatrix} \quad
		T = \begin{bmatrix}
			5 & -1 & \\ 
			-1 & \ddots & \ddots \\
			& \ddots & \ddots & -1\\
			& & -1 & 5
		\end{bmatrix}
	\end{equation*}
	In questo caso i metodi iterativi sono vantaggiosi poiché anche se la matrice è predominante diagonale (e quindi permette Gauss senza problemi), a causa dell'effetto del fill-in lo rende sconveniente.
\end{example}

\begin{example}
	Data una matrice
	\begin{equation*}
		A = \begin{bmatrix}
			n & -1 & \ldots & -1\\ 
			-1 & \ddots & \ddots & \vdots \\
			\vdots & \ddots & \ddots & -1\\
			-1 & \ldots & -1 & n
		\end{bmatrix}
	\end{equation*}
	calcoliamo $P$
	\begin{equation*}
		N = \begin{bmatrix}
			0 & 1 & \ldots & 1\\ 
			1 & \ddots & \ddots & \vdots \\
			\vdots & \ddots & \ddots & 1\\
			1 & \ldots & 1 & 0
		\end{bmatrix} \quad
		P = \begin{bmatrix}
			0 & \frac{1}{n} & \ldots & \frac{1}{n}\\ 
			\frac{1}{n} & \ddots & \ddots & \vdots \\
			\vdots & \ddots & \ddots & \frac{1}{n}\\
			\frac{1}{n} & \ldots & \frac{1}{n} & 0
		\end{bmatrix}
	\end{equation*}
	La sua norma vale
	\begin{equation*}
		\lvert\lvert P \rvert\rvert_{\infty} = \frac{n-1}{n}
	\end{equation*} e quindi converge.
\end{example}

\begin{note}
	Se la norma vale $1$ non posso dire nulla sulla convergenza.
\end{note}

\begin{theorem}
	\begin{equation}
		x^{(k+1)} = Px^{(k)}+q \text{ è convergente } \Leftrightarrow \phi(P)<1
	\end{equation}
\end{theorem}

\begin{example}
		Data una matrice
	\begin{equation*}
		A = \begin{bmatrix}
			1 & & & \alpha\\ 
			-1 & \ddots &  & \\
			& \ddots & \ddots &\\
			& & -1 & 1
		\end{bmatrix}
	\end{equation*}
	troviamo le matrici per la scomposizione
	\begin{equation*}
		M=I \quad
		N = \begin{bmatrix}
			0 & &  & -\alpha\\ 
			1& \ddots & &\\
			 & \ddots & \ddots &\\
			&& 1 & 0
		\end{bmatrix}
	\end{equation*}
	Per quali $\alpha$ il metodo è convergente? \\
	Troviamo la fattorizzazione LU della matrice per il calcolo del determinante e otteniamo così il polinomio caratteristico:
	\begin{equation*}
		x^n + \alpha
	\end{equation*}
	Dobbiamo poi trovare il modulo degli autovalori per poterli confrontare:
	\begin{equation*}
		\begin{split}
			x^n=-\alpha \\
			\lvert\lambda\rvert^n = \lvert -\alpha \rvert \Rightarrow \lvert \lambda \rvert = \sqrt[n]{\lvert \alpha \rvert} \\
			\sqrt[n]{\lvert \alpha \rvert} < 1 \Leftrightarrow \lvert \alpha \rvert < 1
		\end{split}
	\end{equation*}
\end{example}
% !TeX spellcheck = it_IT
\newpage
\section{Equazioni non lineari}
Stiamo considerando equazioni del tipo $f(x)=0$ dove la funzione $f$ non è lineare (quindi non è una retta). Di fronte a questo tipo di equazioni, ci sono due difficoltà:
\begin{itemize}
	\item Non c'è una teoria generale sul \textit{numero} e sull'\textit{esistenza} delle \textbf{soluzioni}
	\item Non esistono metodi diretti di risoluzione
\end{itemize}

\begin{example}
	Determinare il numero di soluzioni reali dell'equazione
	\begin{equation*}
		f(x)=x \log x -1 = 0
	\end{equation*}
	Il primo passo è tracciare un grafico approssimativo di questa funzione:
	\begin{itemize}
		\item \textbf{Dominio}: $x>0$
		\item \textbf{Limiti}:
		\begin{align*}
			& \lim_{x \to + \infty} x \log x -1 = + \infty \\
			& \lim_{x \to 0^+} x \log x = \lim_{x \to 0^+} \frac{log x}{\frac{1}{x}} = \lim_{x \to 0^+} \frac{\frac{1}{x}}{-\frac{1}{x^2}} = \lim_{x \to 0^+} - \frac{x^2}{x} = 0 \Longrightarrow \lim_{x \to 0^+} x log x -1 = -1 
		\end{align*}
		\item \textbf{Derivata prima}: \begin{align*}
			& f'(x) = \log x + x \cdot \frac{1}{x} = \log x + 1 \\
			& f'(x) \geq 0 \Leftrightarrow log x + 1 \geq 0 \Leftrightarrow log x \geq -1 \Leftrightarrow x \geq \frac{1}{e}
		\end{align*}
		\item \textbf{Derivata seconda}: $f''(x)=\frac{1}{x} \geq 0 \quad \forall x > 0$
	\end{itemize}
	\begin{center}
		\begin{tikzpicture}
			\begin{axis}[grid=both,
				xmax=10,ymax=10,
				axis lines=middle,
				samples=100,
				enlargelimits]
				\addplot[green, ultra thick]  {x*ln(x)-1};
			\end{axis}
		\end{tikzpicture}
	\end{center}
	Quindi possiamo dire che 
	\begin{equation*}
		\exists ! \alpha \in \mathbb{R} \vert f(\alpha) = 0
	\end{equation*}
	Ci serve dare un \textbf{intervallo di localizzazione} della soluzione, ad esempio:
	\begin{align*}
		& f(1) = -1 \\
		& f(2) = 2 \log 2 -1 = log 4 -1 \\
		& \Rightarrow \alpha \in [1,2]
	\end{align*}
\end{example}

\subsection{Tecnica della separazione}
A partire da una funzione complessa, mi riconduco a funzioni più semplici e vedo dove si intercettano.
\begin{example}
	\label{example:non_linear_eq}
	Supponiamo di avere
	\begin{equation*}
		x \log x -1 = 0 \Leftrightarrow x \log x = 1 \Leftrightarrow \log x = \frac{1}{x}
	\end{equation*}
	Che sul grafico sono
	\begin{center}
		\begin{tikzpicture}
			\begin{axis}[grid=both,
				xmax=3,ymax=3,
				axis lines=middle,
				restrict x to domain=0:4,
				restrict y to domain=-2:4,
				samples=100,
				enlargelimits]
				\addplot[green, ultra thick]  {x*ln(x)};
				\addplot [blue, ultra thick] {1/x};
			\end{axis}
		\end{tikzpicture}
	\end{center}
\end{example}

\begin{example}
	Data la seguente funzione
	\begin{equation*}
		f(x) = e^x - 2x = 0 \Leftrightarrow e^x = 2x
	\end{equation*}
	È difficile usare il metodo della separazione perché l'intersezione non è facile da trovare. \\
	Usiamo quindi la soluzione grafica:
	\begin{itemize}
		\item \textbf{Dominio}: $\forall x \in \mathbb{R}$ che possiamo scrivere anche come $C^\infty (\mathbb{R})$
		\item \textbf{Limiti}:
		\begin{align*}
			& \lim_{x \to + \infty} e^x -2x = \lim_{x \to + \infty} e^x (1 - \frac{2x}{e^x}) = 0\\
			& \lim_{x \to - \infty} e^x -2x = + \infty\\
		\end{align*}
		\item \textbf{Derivata prima}: \begin{align*}
			& f'(x) = e^x -2 \\
			& f'(x) \geq 0 \Leftrightarrow e^x \geq 2 \Leftrightarrow x \geq \log 2
		\end{align*}
		\item \textbf{Derivata seconda}: $f''(x)=e^x$
	\end{itemize}
	\begin{center}
		\begin{tikzpicture}
			\begin{axis}[grid=both,
				xmax=5,ymax=5,
				axis lines=middle,
				enlargelimits]
				\addplot[green, ultra thick]  {pow(e,x)-2};
			\end{axis}
		\end{tikzpicture}
	\end{center}	
	Calcoliamo adesso il valore in $\log 2$:
	\begin{equation*}
		f(\log 2) = e^{\log 2} - 2 \log 2 = 2-2 \log 2 = 2 (1-\log 2)
	\end{equation*}
\end{example}

\begin{example}
	Data la funzione:
	\begin{equation}
		f(x)= x^3 -6x +1 = 0
	\end{equation}
	In questo esempio abbiamo un polinomio e abbiamo quindi un'\textbf{equazione algebrica} di terzo grado. Quindi sappiamo il numero di soluzioni \textit{complesse}, nel nostro caso $3$. A noi però interessa il numero di soluzioni reali. Sapendo che quelle complesse devono andare sempre in coppia, potremo avere o due soluzioni complesse e una reale oppure tre soluzioni reali. Studiamo la funzione:
	\begin{itemize}
		\item \textbf{Dominio}: $\forall x \in \mathbb{R}$
		\item \textbf{Limiti}:
		\begin{align*}
			& \lim_{x \to + \infty} x^3 -6x +1 = + \infty\\
			& \lim_{x \to - \infty} x^3 -6x +1 = - \infty\\
		\end{align*}
		Sappiamo quindi che esiste sicuramente almeno un punto in cui la funzione vale $0$ per il teorema dell'esistenza degli zeri.
		\item \textbf{Derivata prima}: \begin{align*}
			& f'(x) = 3x^2-6 \\
			& f'(x) = 0 \Leftrightarrow x^2 = 2 \Leftrightarrow x = \pm \sqrt{2} \to x < -\sqrt{2} \land x > \sqrt{2}
		\end{align*}
		\item \textbf{Derivata seconda}: $f''(x)=6x \to f''(x) > 0 \Leftrightarrow x>0$
	\end{itemize}
	\begin{center}
		\begin{tikzpicture}
			\begin{axis}[grid=both,
				xmax=6,ymax=6,
				axis lines=middle,
				restrict x to domain=-4:4,
				restrict y to domain=-10:10,
				samples=100,
				enlargelimits]
				\addplot[blue, ultra thick] (x,x^3-6*x+1);
			\end{axis}
		\end{tikzpicture}
	\end{center}
	Quindi ci sono tre soluzioni reali, che possiamo localizzare come:
	\begin{align*}
		& \beta \in [0, \sqrt{2}]\\
		& \gamma \in [\sqrt{2}, 3]\\
		& \alpha \in [-3,-2]
	\end{align*}
\end{example}
\newpage
\subsection{Metodo di bisezione}
Vogliamo risolvere un'equazione $f(x)=0 \quad f:[a,b]\to \mathbb{R}$ che assumiamo essere \textit{continua} ($f \in C^0([a,b])$). Supponiamo poi che $f(a)f(b) < 0$, quindi la funzione assume valori discordi agli estremi. Questo implica, grazie al teorema di esistenza degli zeri, che $\exists \alpha \in [a,b] \:\vert f(\alpha)=0$, ovvero che esiste almeno un punto di intersezione con l'asse delle x.\\
Assumiamo che esista unico il punto di intersezione:
\begin{equation*}
	\exists ! \alpha \in [a,b] \:\vert\: f(x)=0
\end{equation*}
In questa situazione possiamo costruire diverse successioni che convergono ad $\alpha$ tramite il metodo di bisezione.
\begin{lstlisting}[language=MatLAB]
	a0 = a; b0=b; ya=f(a); yb=f(b);
	for k=1 : inf
		Ck = (a[k-1] + b[k-1])/2; % Calcolo il punto di mezzo dell'intervallo
		y = f(Ck);
		if(y * ya <= 0)
			ak = a[k-1]; bk=Ck; yb=y; % Qui la funzione e' discorde, mi sposto su questo intervallo
		else
			ak=Ck; bk=b[k-1]; ya=y; % Qui e' concorde, mi sposto sull'altro intervallo
		end
	end
\end{lstlisting}
Praticamente prendo il punto medio e controllo quale parte della funzione è concorde o discorde e mi sposto di conseguenza.
\begin{observation}
	Possiamo fare le seguenti osservazioni:
	\begin{enumerate}
		\item $a_k \leq b_k \quad \forall k$
		\item $\alpha \in [a_k, b_k] \quad \forall k$
		\item $b_k - a_k = \frac{b_{k-1}-a_{k-1}}{2} = \ldots = \frac{b_0 - a_0}{2^k}$
	\end{enumerate}
	che implicano:
	\begin{align*}
		& 0 \leq \lvert \alpha - a_k \rvert \leq b_k - a_k = \frac{b_0 - a_0}{2^k}\\
		& 0 \leq \lvert \alpha - b_k \rvert \leq b_k - a_k = \frac{b_0 - a_0}{2^k}\\
		& 0 \leq \lvert C_k  - \alpha \rvert \leq b_{k-1} - a_{k-1}\leq= \frac{b_0 - a_0}{2^k}
	\end{align*}
	e portando $k$ all'infinito, $a_k$, $b_k$ e $C_k$ tendono tutti ad $\alpha$.
\end{observation}

\begin{example}
	Supponiamo di voler risolvere
	\begin{equation*}
		f(x)=x-\cos x = 0
	\end{equation*}
	\newpage
	Proviamo prima con la \textbf{tecnica della separazione}:
	\begin{center}
		\begin{tikzpicture}
			\begin{axis}[grid=both,
				xmax=3,ymax=3,
				axis lines=middle,
				restrict x to domain=-5:5,
				restrict y to domain=-5:5,
				samples=100,
				enlargelimits]
				\addplot[green, ultra thick]  {cos(deg(x))};
				\addplot [blue, ultra thick] {x};
			\end{axis}
		\end{tikzpicture}
	\end{center}
	Osservando  attentamente notiamo che
	\begin{equation*}
		\exists ! \alpha \in [0,1]
	\end{equation*}
	poiché fuori da questo intervallo la funzione $x$ assume valori maggiori di $1$ o minori di $-1$ che non sono valori che può assumere la funzione $cos(x)$.\\
	Per poter usare il \textbf{metodo della bisezione} devo verificare che $f(x)$ sia continua e la sua monotonia.\\
	Come estremi per il metodo possiamo porre $[a,b]=[0,1]$.
	\begin{equation*}
		\lvert C_k - \alpha \rvert \leq \frac{b_0 - a_0}{2^k} \leq \epsilon
	\end{equation*}
	Posso a priori decidere la distanza dal mio numero finale facendo $\log_2 \frac{1}{\epsilon}$ passi.
\end{example}

\subsubsection{Approssimazione della radice}
Supponiamo di voler risolvere
\begin{equation*}
	x^2-2=0
\end{equation*}
Sappiamo che le soluzioni sono $x=\pm \sqrt{2}$ e che il grafico è:
\begin{center}
	\begin{tikzpicture}
		\begin{axis}[grid=both,
			xmax=3,ymax=3,
			axis lines=middle,
			restrict x to domain=-5:5,
			restrict y to domain=-5:5,
			samples=100,
			enlargelimits]
			\addplot[green, ultra thick]  {x^2-2};
		\end{axis}
	\end{tikzpicture}
\end{center}
Il calcolatore non sa come calcolare $\sqrt{2}$, il modo migliore per farglielo fare è approssimare questa funzione con il metodo di bisezione nell'intervallo $[0,2]$.

\subsubsection{Svantaggi}
Questo metodo ha fondamentalmente due svantaggi:
\begin{itemize}
	\item Ha difficoltà ad essere esteso ai numeri \textbf{complessi}, ma per il nostro corso non ci riguarda
	\item Il \textbf{costo computazionale} è molto influenzato dal numero di valutazioni della funzioni che faccio e, dato che non conosco precisamente quali calcoli dovrà fare, potrebbe essere un problema. Nel metodo di bisezioni faccio una valutazione per passo e il numero di passi potrebbe essere elevato.
\end{itemize} 

\subsection{Metodo del punto fisso}
Questo metodo prevede di riscrivere il nostro problema come:
\begin{equation*}
	f(x) \Longleftrightarrow x = g(x)
\end{equation*}
\begin{example}
	Alcuni esempi:
	\begin{align*}
		&f(x)=0 \Longleftrightarrow x = x-f(x) \quad g(x) = x-f(x)\\
		& f(x)=0 \Longleftrightarrow x=x-\frac{f(x)}{h(x)} \quad g(x)=x-\frac{f(x)}{h(x)}
	\end{align*}
\end{example}
Poi si genera una successione da un punto iniziale:
\begin{equation*}
	\begin{cases}
		x_0 \in \mathbb{R} \\
		x_{\alpha +1} = g(x_\alpha)
	\end{cases}
\end{equation*}

\begin{example}
	\label{example:fixed_point}
	Partiamo dalla seguente equazione:
	\begin{equation*}
		x - \sqrt{x} = 0 \quad x \geq 0
	\end{equation*}
	Che posso riscrivere come:
	\begin{equation*}
		x = \sqrt{x} \Longleftrightarrow x^2 = x \Longleftrightarrow x = 0 \lor x=1
	\end{equation*}
	Notiamo come abbiamo ad esempio due possibili modi per scrivere l'equazione:
	\begin{align*}
		& x- \sqrt{x} = 0 \Longleftrightarrow x = \sqrt{x} \quad g_1(x)= \sqrt{x} \\
		& x- \sqrt{x} = 0 \Longleftrightarrow x^2 = x \quad g_2(x)=x^2
	\end{align*}
	Vediamo la successione generata da $g_1(x)$:
	\begin{equation*}
		\begin{cases}
			x_0 \in \mathbb{R}^+ \\
			x_{\alpha +1} = g_1(x_\alpha) = \sqrt{x_\alpha}
		\end{cases}
	\end{equation*}
	Che graficamente è
	\begin{center}
		\begin{tikzpicture}
			\begin{axis}[grid=both,
				xmax=2,ymax=2,
				axis lines=middle,
				restrict x to domain=0:2,
				restrict y to domain=0:2,
				samples=100,
				enlargelimits]
				\addplot[green, ultra thick]  {sqrt(x)};
				\addplot[blue, ultra thick]  {x};
			\end{axis}
		\end{tikzpicture}
	\end{center}
	Con questa iterazione, se faccio i vari calcoli, da qualunque punto iniziale io scelga finirò alla soluzione $1$ (intersezione di $x$ e $\sqrt{x}$).\\
	Se scelgo invece la successione generata da $g_2(x)$, ottengo l'effetto contrario e trovo solo la soluzione $0$ quando parto da $0 < x_0 < 1 \lor -1 < x_0 < 0$, altrimenti la successione diverge a $+\infty$ per $x_0>1 \lor x_0 < -1$.
	\begin{center}
		\begin{tikzpicture}
			\begin{axis}[grid=both,
				xmax=2,ymax=2,
				axis lines=middle,
				restrict x to domain=-2:3,
				restrict y to domain=-1:3,
				samples=100,
				enlargelimits]
				\addplot[green, ultra thick]  {x^2};
				\addplot[blue, ultra thick]  {x};
			\end{axis}
		\end{tikzpicture}
	\end{center}
\end{example}

\begin{theorem}[Teorema del punto fisso]
	Sia $g \in C^1([a,b])$, $\alpha \in (a,b)$ e $g(\alpha) = \alpha$ un \textbf{punto fisso}.\\
	\textit{Ipotesi}: se $\exists \rho >0 \:\vert\: \forall x \in [\alpha - \rho \alpha + \rho]=I_\alpha \geq [a,b]$ si ha che $\lvert g(x) \rvert < 1$.\\
	\textit{Tesi}: allora $\forall x_0 \in I_\alpha$ la successione generata dal metodo $x_{k+1} = g(x_k)$ soddisfa le seguenti proprietà:
	\begin{itemize}
		\item $x_= \in I_\alpha \quad \forall k$
		\item $x_k \to \alpha$
	\end{itemize}
\end{theorem}

\begin{observation}
	È importante nel precedente teorema che l'intervallo $[a,b]$ sia centrato per garantire che da qualunque punto si inizi si riduca la distanza dal punto fisso senza però uscire dall'intervallo.
\end{observation}

\begin{demostration}
	Sappiamo che se $g(x)$ è continua allora anche $\lvert g'(x) \rvert$ è continua. Per il \textbf{teorema di Weirstrass} sappiamo che una funzione continua su un intervallo limitato ammette \textit{massimo}, nel nostro caso $\lambda < 1$. \\
	Dimostriamo per induzione che $x_0 \in I_\alpha \Rightarrow \lvert x_k -\alpha \rvert \leq \lambda^k \rho \quad \forall k \geq 0$:
	\begin{itemize}
		\item \textit{Passo base}: per $k=0$ abbiamo che: 
		\begin{equation*}
			\lvert x_0 - \alpha \rvert \leq \lambda^0 \rho = \alpha
		\end{equation*}
		che è vero per le ipotesi: $x_0 \in I_\alpha$.
		\item \textit{Passo induttivo}: posso scrivere:
		\begin{equation*}
			\lvert x_{k+1} - \alpha \rvert = \lvert g(x_k) - g(\alpha) \rvert = \lvert g'(\psi k)(x_k - \alpha) \rvert = \lvert g'(\psi k) \rvert  \lvert x_k - \alpha \rvert \leq \lambda \lambda^k \rho = \lambda ^ {k+1} \rho
		\end{equation*}
	\end{itemize}
\end{demostration}

\begin{example}
	Partiamo dalla funzione dell'esempio \ref{example:fixed_point}, in particolare da $g_1$:
	\begin{equation*}
		g'_1(x) = \frac{1}{2 \sqrt{x}} \Longrightarrow \lvert g'_1(x) \rvert = \lvert \frac{1}{2} \frac{1}{\sqrt{x}} \rvert = \frac{1}{2 \sqrt{x}}
	\end{equation*}
	Vediamo dove il suo modulo è minore di $1$:
	\begin{equation*}
		\frac{1}{2 \sqrt{x}} < 1 \Leftrightarrow 2 \sqrt{x} > 1 \Leftrightarrow \sqrt{x} > \frac{1}{2} \Leftrightarrow x> \frac{1}{4}
	\end{equation*}
	Quindi comunque si scelga un intervallo centrato in $1$ tale che $x>\frac{1}{4}$ (al massimo potrò avere $[\frac{1}{4}, \frac{5}{4}]$) allora la successione converge.\\
	Vediamo la funzione $g_2$:
	\begin{equation*}
		g'_2(x) = 2x \Longrightarrow \lvert g'_2(x) \rvert = \lvert 2x \rvert = 2 \lvert x \rvert
	\end{equation*}
	Vediamo dove il suo modulo è minore di $1$:
	\begin{equation*}
		2 \lvert x \rvert < 1 \Leftrightarrow \lvert x \rvert < \frac{1}{2}
	\end{equation*}
	Quindi comunque si scelga un intervallo centrato in $0$ tale che $x$ è compreso tra $-\frac{1}{2}$ e $\frac{1}{2}$, allora la successione converge.
\end{example}

\begin{theorem}
	Prendiamo $g \in C^1([a,b]) \quad g(\alpha)=\alpha \quad \alpha \in (a,b)$. \\
	Se $\lvert g'(\alpha) \rvert < 1$, prendiamo $h(x) = \lvert g'(x) \rvert -1$ quindi:
	\begin{equation*}
		h(\alpha) = \lvert g'(\alpha) \rvert -1 < 0 \Rightarrow \exists [\alpha-\rho, \alpha + \rho] = I_\alpha \:\vert\: \forall x \in I_\alpha \quad h(x)<0 \Leftrightarrow \lvert g'(x) \rvert <1
	\end{equation*}
	Quindi posso dire che $\alpha$ è \textbf{punto attrattivo} (se parto vicino ad $\alpha$ convergo su di esso, esiste un intorno e c'è convergenza locale). Altrimenti è \textbf{punto repulsivo}. Se è uguale, allora è \textbf{neutrale}.
\end{theorem}

\begin{example}
	Prendiamo la funzione
	\begin{equation*}
		f(x) = e^{-x} - x = 0
	\end{equation*}
	\begin{enumerate}
		\item Utilizziamo il metodo di separazione grafica per determinare che esiste unica la soluzione:
		\begin{equation*}
			e^{-x} - x = 0 \Leftrightarrow e^{-x} = x
		\end{equation*}
		\begin{center}
			\begin{tikzpicture}
				\begin{axis}[grid=both,
					xmax=2,ymax=2,
					axis lines=middle,
					restrict x to domain=-2:3,
					restrict y to domain=-1:3,
					samples=100,
					enlargelimits]
					\addplot[green, ultra thick]  {e^(-x)};
					\addplot[blue, ultra thick]  {x};
				\end{axis}
			\end{tikzpicture}
		\end{center}
		E ho trovato un intervallo di separazione $\alpha \in [0, 1]$.
		\item Consideriamo il metodo iterativo
		\begin{equation*}
			x_{k+1} = e^{-x_k} \quad k \geq 0
		\end{equation*}
		\begin{itemize}
			\item Dire se l'intervallo è localmente convergente. Sappiamo che 
			\begin{equation*}
				g'(x)=-e^{-x} \Rightarrow \lvert g'(x) \rvert = e^{-x}
			\end{equation*}
			quindi
			\begin{equation*}
				e^{-x} < 1 \Leftrightarrow -x < 0 \Leftrightarrow x >0
			\end{equation*}
			e posso dire che $\lvert g'(\alpha)  \rvert< 1$ e quindi il metodo è \textbf{localmente convergente}.
			\item Si dica come scegliere un punto iniziale a partire dal quale il metodo converge. Per localizzare $\alpha$, utilizzo il metodo di bisezione:
			\begin{equation*}
				x=\frac{1}{2} \Rightarrow g(\frac{1}{2}) = \frac{1}{\sqrt{e}} > \frac{1}{2}
			\end{equation*}
			Quindi posso scegliere $1$ perché ho abbastanza spazio dall'altro lato per centrare l'intervallo in $\alpha$.
		\end{itemize}
	\end{enumerate}
\end{example}
\newpage
\begin{example}
	Prendiamo la funzione dell'esempio \ref{example:non_linear_eq} e vediamo che
	\begin{equation*}
		\alpha \in [1,2]
	\end{equation*}
	Consideriamo due metodi iterativi:
	\begin{itemize}
		\item
		\begin{equation*}
			x_{k+1} = \frac{1}{\ln {x_k}} \quad k \geq 0
		\end{equation*}
		\item 
		\begin{equation*}
			x_{k+1} = e^{\frac{1}{x_k}} \quad k \geq 0
		\end{equation*}
	\end{itemize}
	Verifichiamo le derivate:
	\begin{itemize}
		\item 
		\begin{equation*}
			g'(x) = \frac{- \frac{1}{x}}{\ln ^ 2 {x}} \Rightarrow \lvert \frac{- \frac{1}{x}}{\ln ^ 2 {x}} \rvert  = \frac{1}{x \ln^2 x}
		\end{equation*}
		La valuto in $\alpha$ tenendo a mente che è punto fisso ($\alpha = \frac{1}{\alpha}$):
		\begin{equation*}
			\lvert g'(\alpha)= \frac{1}{\alpha \ln ^2 \alpha} = \frac{\alpha^2}{\alpha} = \alpha > 1
		\end{equation*}
		E quindi per questo metodo è un \textbf{punto repulsivo}.
		\item \begin{equation*}
			g'(x) = e^\frac{1}{x}(-\frac{1}{x^2})
		\end{equation*}
		La valuto in $\alpha$:
		\begin{equation*}
			\lvert g'(\alpha)= e^\frac{1}{\alpha}(-\frac{1}{\alpha^2}) = \frac{1}{\alpha} < 1
		\end{equation*}
		E quindi per questo metodo è un \textbf{punto attrattivo}.
	\end{itemize}
\end{example}
% !TeX spellcheck = it_IT
\newpage
\section{Teoremi}
\subsection{Maggiorazione dell'errore di macchina}
\begin{theorem}
	Sia
	\begin{equation*}
		x \in \mathbb{R}
	\end{equation*}
	con
	\begin{equation*}
		\omega \leq \lvert x \rvert \leq \Omega
	\end{equation*}
	vale
	\begin{equation}
		\lvert \epsilon_x \rvert = \bigg\lvert \frac{\text{trn}(x) - x}{x} \bigg\rvert \leq u = B^{1-t}
	\end{equation}
\end{theorem}

\subsubsection{Dimostrazione}
Definiamo $x$ come
\begin{equation*}
	x = (-1)^s B^p \alpha
\end{equation*}
L'errore assoluto $\lvert \text{trn}(x) - x \rvert$ è maggiorato dalla distanza tra due numeri di macchina consecutivi e quindi:
\begin{equation*}
	\lvert \text{trn}(x) - x \rvert \leq B^{p-t}
\end{equation*}
Inoltre vale
\begin{equation*}
	\lvert x \rvert \geq B^{p-1}
\end{equation*}
Quindi:
\begin{equation*}
		\lvert \epsilon_x \rvert = \bigg\lvert \frac{\text{trn}(x) - x}{x} \bigg\rvert \leq \frac{B^{p-t}}{B^{p-1}} = B^{1-t} = u
\end{equation*}

\subsection{Errore analitico e totale}
\begin{definition}[Errore analitico]
	Nel calcolo di $h(x) \neq 0$ tramite la sua approssimazione $f(x)$, l'errore analitico è
	\begin{equation}
		\epsilon_{an} = \frac{f(x)-h(x)}{h(x)}
	\end{equation}
\end{definition}
\begin{definition}[Errore totale]
	Combinando l'errore analitico con quello \textbf{algoritmo} e \textbf{inerente} otteniamo:
	\begin{equation}
		\epsilon_{tot} = \frac{g(\tilde{x}) - h(x)}{h(x)} \doteq \epsilon_{an}  + (\epsilon_{in} + \epsilon_{alg})
	\end{equation}
\end{definition}

\newpage
\subsection{Localizzazione degli autovalori}
\begin{definition}[Cerchi di Gershgorin]
	Sia
	\begin{equation*}
		A = (a_{i,j}) \in \mathbb{C}^{n \times n}
	\end{equation*}
	Definiamo i cerchi di Gershgorin $K_i$ con $1 \leq i \leq n$ come
	\begin{equation}
		K_i = \{z \in \mathbb{C} : \lvert z - a_{i,i} \rvert \leq \sum_{j=i, j \neq i}^{n}\lvert a_{i,j} \rvert\}
	\end{equation}
\end{definition}

\begin{theorem}
	Vale
	\begin{equation}
		\lambda \text{ autovalore di } A \Longrightarrow \lambda \in \bigcup_{i=1}^n K_i
	\end{equation}
\end{theorem}

\subsubsection{Dimostrazione}
Sia $\lambda$ \textbf{autovalore} di $A$ con corrispondente \textbf{autovettore destro} $\mathbf{x}$.\\
La relazione $A\mathbf{x} = \lambda \mathbf{x}$ implica che
\begin{equation*}
	\sum_{j=1}^{n}a_{i,j}x_j = \lambda x_i \quad\quad 1 \leq i \leq n
\end{equation*}
da cui
\begin{equation}
	(\lambda - a_{i,i})x_i = \sum_{j=1,j \neq i}^{n}a_{i,j}x_j \quad\quad 1 \leq i \leq n
\end{equation}
Sia $p$ l'indice di una componente di modulo massimo di $\mathbf{x}$, ad esempio $\lvert x_p \rvert = \lvert\lvert \mathbf{x} \rvert\rvert _\infty$. Dato che $\mathbf{x} \neq 0$, allora $\lvert x \rvert >0$.\\
Poniamo $i=p$ nell'equazione precedente e otteniamo
\begin{equation*}
	(\lambda - a_{p,p})x_p = \sum_{j=1,j \neq p}^{n}a_{p,j}x_j
\end{equation*}
che, con i valori assoluti, diventa
\begin{equation*}
	\lvert(\lambda - a_{p,p})x_p\rvert = \lvert \lambda - a_{p,p} \rvert \lvert x_{p} \rvert= \bigg\lvert\sum_{j=1,j \neq p}^{n}a_{p,j}x_j \bigg\rvert \leq \sum_{j=1,j \neq p}^{n}\lvert a_{p,j}\rvert \lvert x_j\rvert
\end{equation*}
Se dividiamo entrambi i membri per $\lvert x_p \rvert$ otteniamo
\begin{equation*}
	\lvert \lambda - a_{p,p} \rvert \leq \sum_{j=1,j \neq p}^{n}\lvert a_{p,j}\rvert \frac{\lvert x_j\rvert}{\lvert x_p \rvert} \leq \sum_{j=1,j \neq p}^{n}\lvert a_{p,j}\rvert
\end{equation*}
E questo implica che $\lambda \in K_p$.

\newpage
\subsection{Esistenza della fattorizzazione LU}
\begin{theorem}
	Sia
	\begin{equation*}
		A \in \mathbb{R}^{n \times n}
	\end{equation*}
	Se $A(1:k,1:k)$ è \textbf{invertibile} per $k=1,2,\ldots, n-1$, allora esiste unica la fattorizzazione LU di $A$.
\end{theorem}
\subsubsection{Dimostrazione}
Per \textbf{induzione} sulla dimensione $n$.
\begin{itemize}
	\item \textbf{Caso base}: per $n=1$ vale che $A=[a] = [1][a]$ che è l'unica fattorizzazione LU
	\item \textbf{Ipotesi induttiva}: supponiamo che il teorema sia vero per matrici di ordine $m \leq n-1$
	\item \textbf{Passo induttivo}: la fattorizzazione può essere scritta come
	\begin{equation*}
		\begin{bmatrix}
			A(1:n-1, 1:n-1) & \mathbf{z} \\
			\mathbf{v}^T & \alpha
		\end{bmatrix} = \begin{bmatrix}
		L(1:n-1, 1:n-1) & 0 \\
		\mathbf{\omega}^T & 1
		\end{bmatrix} \cdot \begin{bmatrix}
		U(1:n-1, 1:n-1) & \mathbf{y} \\
		0^T & \beta
		\end{bmatrix}
	\end{equation*}
	Questo è equivalente al sistema
	\begin{equation*}
		\begin{cases}
			A(1:n-1,1:n-1) = L(1:n-1,1:n-1) \cdot U(1:n-1, 1:n-1) \\
			\mathbf{z} = L(1:n-1, 1:n-1)\mathbf{y} \\
			\mathbf{v}^T=\omega^T \cdot U(1:n-1, 1:n-1) \\
			\alpha = \mathbf{\omega}^T \mathbf{y} + \beta
		\end{cases}
	\end{equation*}
	\begin{enumerate}
		\item Per ipotesi induttiva $A(1:n-1, 1:n-1)$ è invertibile e di conseguenza $L(1:n-1,1:n-1)$ e $U(1:n-1, 1:n-1)$ sono i suoi fattori triangolari 
		\item  Per ipotesi induttiva $A(1:n-1, 1:n-1)$ è invertibile, quindi lo è anche $U(1:n-1, 1:n-1)$ lo è, di conseguenza le equazioni 2 e 3 ammettono soluzione unica
		\item Dati $\mathbf{\omega}$ e $\mathbf{y}$ l'equazione 4 permette di identificare univocamente $\alpha$ e $\beta$
	\end{enumerate}
\end{itemize}

\subsection{Convergenza con norma}
\begin{theorem}
	Un metodo iterativo è \textbf{convergente} se esiste una norma matriciale indotta da una norma vettoriale $\lvert\lvert \cdot\rvert\rvert$ su $\mathbb{R}^n$ tale per cui
	\begin{equation*}
		\lvert\lvert P \rvert\rvert < 1
	\end{equation*}
\end{theorem}

\subsubsection{Dimostrazione}
Dalle relazioni
\begin{equation*}
	\mathbf{x}^{(k+1)} = P\mathbf{x}^{(k)}+\mathbf{q} \quad\quad\quad \mathbf{x}=P\mathbf{x}+q
\end{equation*}
segue
\begin{equation*}
	\mathbf{e}^{(k+1)} = \mathbf{x}^{(k+1)} - \mathbf{x} = P(\mathbf{x}^{(k)}-x)=P\mathbf{e}^{(k)} \quad\quad k \geq 0
\end{equation*}
e quindi
\begin{equation*}
	\mathbf{e}^{(k+1)} = P^{k+1}e^{(0)} \quad\quad k\geq 0
\end{equation*}
Passando alla norma vettoriale
\begin{equation*}
	\lvert \lvert \mathbf{e}^{(k+1)} \rvert\rvert= \lvert\lvert P^{k+1}e^{(0)}\rvert\rvert  \leq  \lvert\lvert P^{k+1}\rvert\rvert\: \lvert\lvert e^{(0)}\rvert\rvert  
\end{equation*}
da cui
\begin{equation*}
	0 \leq \lvert \lvert \mathbf{e}^{(k+1)} \rvert\rvert \leq \lvert\lvert P\rvert\rvert^{k+1}\lvert\lvert e^{(0)}\rvert\rvert  
\end{equation*}
da cui per il \textbf{teorema del confronto}
\begin{equation*}
	\lim_{k \to +\infty} \lvert\lvert \mathbf{e}^{(k+1)} \rvert\rvert = \lim_{k \to +\infty} \lvert\lvert \mathbf{x}^{(k+1)}-\mathbf{x} \rvert\rvert = 0 \quad\quad \forall\mathbf{e}^{(0)},\mathbf{x}^{(0)}
\end{equation*}

\subsection{Convergenza con raggio spettrale}
\begin{theorem}
	Se il metodo iterativo è convergente allora $\rho(P)<1$.
\end{theorem}
\subsubsection{Dimostrazione}
Sia $\lambda$ tale che $\lvert \lambda \rvert = \rho(P)$ e $\mathbf{v}$ il corrispondente \textbf{autovettore}.\\
Sia $\mathbf{x}^{(0)} = \mathbf{x} + \mathbf{v}$ con $\mathbf{x} = A^{-1} \mathbf{b}$ soluzione del sistema $A\mathbf{x} = \mathbf{b}$.\\
La successione con punto iniziale $\mathbf{x}^{(0)}$ converge a $\mathbf{x}$.\\
Si ha
\begin{equation*}
	\mathbf{e}^{(k+1)} = P^{k+1}\mathbf{e}^{(0)} = P^{k+1}\mathbf{v} = \lambda^{k+1}\mathbf{v}
\end{equation*}
da cui
\begin{equation*}
	\lvert\lvert \mathbf{e}^{(k+1)} \rvert\rvert = \lvert\lvert \lambda^{k+1}\mathbf{v} \rvert\rvert = \lvert \lambda\rvert^{k+1} \lvert\lvert \mathbf{v} \rvert\rvert
\end{equation*}
e quindi
\begin{equation*}
	\lim_{k \to +\infty} \lvert \lambda \rvert ^k = 0
\end{equation*}
che implica
\begin{equation*}
	\lvert \lambda \rvert < 1
\end{equation*}

\subsection{Predominanza diagonale}
\begin{theorem}
	Se $A = (a_{i,j}) \in \mathbb{R}^{n \times n}$ è \textbf{predominante diagonale} allora valgono:
	\begin{enumerate}
		\item $A$ è \textbf{invertibile}
		\item I metodi di Jacobi e Gauss-Seidel sono applicabili
		\item I metodi di Jacobi e Gauss-Seidel sono convergenti
	\end{enumerate}
\end{theorem}

\subsubsection{Dimostrazione}
\begin{enumerate}
	\item Dal teorema di Gerschgorin vale
	\begin{equation*}
		\lvert 0 - a_{i,i} \rvert = \lvert a_{i,i} \rvert > \sum_{j=1, j\neq i}^{n} \lvert a_{i,j} \rvert \quad\quad 1 \leq i \leq n
	\end{equation*}
	e dunque
	\begin{equation*}
		0 \notin \bigcup_{i=1}^n K_i
	\end{equation*}
	\item \begin{equation*}
		\lvert a_{i,i} \rvert > \sum_{j=1, j\neq i}^n \lvert a_{i,j} \rvert \geq 0 \Longrightarrow \lvert a_{i,i} \rvert \neq 0 \quad\quad 1 \leq i \leq n
	\end{equation*}
	\item  Dalla relazione
	\begin{equation*}
		det(P-\lambda I_n) = det(M^{-1}N-\lambda I_n) = det(N - \lambda M) = (-1)^n det(\lambda M - N)
	\end{equation*}
	Quindi $\lambda$ è autovalore di $P$ se e solo se $det(\lambda M - N) = 0$.\\\\
	Assumiamo che $\lvert \lambda \rvert \geq 1$. Vediamo che la matrice $\lambda M - N$ è predominante diagonale:
	\begin{equation*}
		\lvert a_{i,i} \rvert > \sum_{j=1, j\neq i}^n \lvert a_{i,j} \rvert = \sum_{j=1}^{i-1} \lvert a_{i,j} \rvert + \sum_{j=i+1}^{n} \lvert a_{i,j} \rvert \quad\quad 1 \leq i \leq n
	\end{equation*}
	\begin{equation*}
		\lvert \lambda \rvert \lvert a_{i,i} \rvert >\lvert \lambda \rvert  \sum_{j=1}^{i-1} \lvert a_{i,j} \rvert + \lvert \lambda \rvert \sum_{j=i+1}^{n} \lvert a_{i,j} \rvert \geq \lvert \lambda \rvert  \sum_{j=1}^{i-1} \lvert a_{i,j} \rvert + \sum_{j=i+1}^{n} \lvert a_{i,j} \rvert 
	\end{equation*}
	e quindi vale la predominanza diagonale
	\begin{equation*}
		\lvert \lambda a_{i,i} \rvert >\sum_{j=1}^{i-1} \lvert \lambda a_{i,j} \rvert + \sum_{j=i+1}^{n} \lvert a_{i,j} \rvert 
	\end{equation*}
	\begin{equation*}
		\lvert \lambda a_{i,i} \rvert >\sum_{j=1}^{i-1} \lvert a_{i,j} \rvert + \sum_{j=i+1}^{n} \lvert a_{i,j} \rvert 
	\end{equation*}
	Per il punto 1 sappiamo che allora la matrice è anche invertibile:
	\begin{equation*}
		\forall \lambda \in \mathbb{C}, \lvert \lambda \rvert \geq 1 \Longrightarrow det(\lambda M - N) \neq 0
	\end{equation*}
	Quindi per gli autovalori delle matrici di iterazione deve valere 
	\begin{equation*}
		\lvert \lambda \rvert < 1
	\end{equation*}
	e dunque
	\begin{equation*}
		\rho(P) < 1
	\end{equation*}
\end{enumerate}

\subsection{Convergenza del metodo di bisezione}
\begin{theorem}
	Sia $f:[a,b] \to \mathbb{R}$ con:
	\begin{itemize}
		\item $f \in C^0([a,b])$
		\item $f(a)f(b) <0$
	\end{itemize}
	allora si ha, con $f(\epsilon)=0$,
	\begin{equation}
		\lim_{k \to +\infty}a_k = \lim_{k \to +\infty} b_k = \lim_{k \to +\infty} c_k = \epsilon \in [a,b]
	\end{equation}
\end{theorem}
\subsubsection{Dimostrazione}
Per costruzione valgono:
\begin{itemize}
	\item $a_{k+1} \geq a_k$
	\item $b_{k+1} \leq b_k$
	\item $c_k \in [a_k,b_k] \subset [a,b]$
	\item $0 \leq b_k - a_k \leq \frac{b-a}{2^{k-1}}$
	\item $f(a_k)f(b_k) \leq 0$
	\item $k \geq 1$
\end{itemize}
Quindi esistono $\epsilon, \eta \in [a,b]$ tali che
\begin{equation*}
	\lim_{k \to +\infty} a_k = \epsilon \quad\quad\quad \lim_{k \to +\infty}b_k = \eta
\end{equation*}
Per il \textbf{teorema del confronto} segue che
\begin{equation*}
	\epsilon = \lim_{k \to +\infty}a_k = \lim_{k \to +\infty} b_k = \eta = \lim_{k \to +\infty} c_k
\end{equation*}
Per la continuità di $f$ si ha
\begin{equation*}
	\lim_{k \to +\infty} f(a_k)f(b_k) = f(\epsilon)^2 \leq 0
\end{equation*}
che implica
\begin{equation*}
	f(\epsilon)=0
\end{equation*}
\subsection{Limite metodo funzionale}
\begin{theorem}
	Sia $g:[a,b]\to \mathbb{R}$ con:
	\begin{itemize}
		\item $g \in C^1([a,b])$
		\item $g(\epsilon) = \epsilon$
		\item $\epsilon \in (a,b)$
	\end{itemize}
	Se esiste $\rho > 0$ tale che:
	\begin{equation*}
		\lvert g'(x) \rvert < 1 \quad \forall x \in [\epsilon - \rho, \epsilon + \rho] = I_{\epsilon} \subset [a,b]
	\end{equation*}
	allora per ogni $x_o \in I_\epsilon$ la successione generata soddisfa
	\begin{enumerate}
		\item $x_k \in I_\epsilon \quad\quad \forall k \geq 0$
		\item $\lim_{k \to +\infty} x_k = \epsilon$
	\end{enumerate}
\end{theorem}
\subsubsection{Dimostrazione}
Dal \textbf{teorema di Weirstrass}, dato che $g'$ è continua e $I_\epsilon$ chiuso e limitato, vale
\begin{equation*}
	\lambda = \max_{x \in I_\epsilon}\lvert g'(x) \rvert < 1
\end{equation*}
Si dimostra che la successione generata soddisfa:
\begin{equation*}
	\lvert x_k - \epsilon \rvert \leq \lambda^k \rho \quad\quad k \geq 0
\end{equation*}
da cui seguono i due punti:
\begin{enumerate}
	\item $\lvert x_k - \epsilon \rvert \leq \lambda^k \rho \leq \rho \Longrightarrow x_k \in I_\epsilon$
	\item per il \textbf{teorema del confronto} $0 \leq \lvert x_k -\epsilon \rvert \leq \lambda^k \rho \Longrightarrow \lim_{k \to +\infty} \lvert x_k - \epsilon \rvert = 0$
\end{enumerate}
Procediamo per \textbf{induzione} su $k$:
\begin{itemize}
	\item \textbf{Passo base}: per $k=0$ si ha $\lvert x_0 - \epsilon \rvert \leq \lambda^0 \rho = \rho$
	\item \textbf{Ipotesi induttiva}: assumiamo che valga la disequazione precedente fino a $k$
	\item \textbf{Passo induttivo}: per il \textbf{teorema di Lagrange} vale
	\begin{equation*}
		\lvert x_{k+1} - \epsilon \rvert = \lvert g(x_k) - g(\epsilon) \rvert = \lvert g'(\eta_k)(x_k - \epsilon) \rvert = \lvert g'(\eta_k)\rvert \lvert (x_k - \epsilon) \rvert \quad\quad \lvert \eta_k - \epsilon \rvert \leq \lvert x_k - \epsilon \rvert
	\end{equation*} 
	Per ipotesi induttiva segue che $\eta_k \in I_\epsilon$ e quindi
	\begin{equation*}
		\lvert x_{k+1} - \epsilon \rvert = \lvert g(\eta_k)\rvert\lvert(x_k - \epsilon) \rvert \leq \lambda \lambda^k \rho = \lambda^{k+1}\rho
	\end{equation*}
\end{itemize}
\subsection{Convergenza locale metodo funzionale}
\begin{theorem}
	Sia $g:[a,b]\to \mathbb{R}$ con:
	\begin{itemize}
		\item $g \in C^1([a,b])$
		\item $g(\epsilon) = \epsilon$
		\item $\epsilon \in (a,b)$
	\end{itemize}
	Se esiste $\lvert g'(\epsilon) \rvert  < 1 \rvert$ allora il metodo è localmente convergente in $\epsilon$.
\end{theorem}
\subsubsection{Dimostrazione}
Sia $h: [a,b] \to \mathbb{R}$ tale che $h(x) = \lvert g'(x) \rvert -1$. Si ha che $h(x) \in C^0([a,b])$ e $h(\epsilon) < 0$.\\
Dal \textbf{teorema della permanenza del segno} abbiamo che esiste
\begin{equation*}
	I_\epsilon = [\epsilon-\rho, \epsilon+\rho] \subset [a,b]
\end{equation*}
tale che
\begin{equation*}
	h(x) = \lvert g'(x) \rvert -1 < 0 \quad\quad \forall x \in I_\epsilon
\end{equation*}
La tesi segue dal teorema precedente.
\subsection{Convergenza locale metodo delle tangenti}
\begin{theorem}
	Sia $f:[a,b]\to \mathbb{R}$ con:
	\begin{itemize}
		\item $f \in C^2([a,b])$
		\item $f(\epsilon) = \epsilon$
		\item $\epsilon \in (a,b)$
		\item $f'(\epsilon) \neq 0$
	\end{itemize}
	allora il metodo è localmente convergente, ovvero \\
	se esiste $\rho > 0$ tale che per ogni $x_0 \in [\epsilon - \rho, \epsilon + \rho] = I_\epsilon \subset [a,b]$ la successione soddisfa:
	\begin{enumerate}
		\item $x_k \in I_\epsilon \quad\quad \forall k\geq 0$
		\item $\lim_{k \to +\infty} x_k = \epsilon$
	\end{enumerate}
	Se la successione soddisfa $x_k \neq \epsilon \quad k \geq 0$ allora la convergenza è almeno \textbf{quadratica}.
\end{theorem}
\subsubsection{Dimostrazione}
Da $f'(\epsilon) \neq 0$ per il \textbf{teorema della permanenza del segno} segue che esiste
\begin{equation*}
	I'_\epsilon = [\epsilon - \rho', \epsilon + \rho'] \subset [a,b]
\end{equation*}
tale che
\begin{equation*}
	f'(x) \neq 0 \quad\quad \forall x \in I'_\epsilon
\end{equation*}
Si verifica che la funzione di iterazione $g: I'\epsilon \to \mathbb{R}$ definita come
\begin{equation*}
	g(x) = x - \frac{f(x)}{f'(x)}
\end{equation*}
soddisfa
\begin{enumerate}
	\item $g \in C^1(I'_\epsilon)$
	\item $g'(\epsilon) = \frac{f(\epsilon)f''(\epsilon)}{(f'(\epsilon))^2}=0$
\end{enumerate}
e quindi la tesi segue dal teorema precedente.\\
Per la stima della velocità di convergenza, dallo \textbf{sviluppo di Taylor} al secondo ordine otteniamo
\begin{equation*}
	0 = f(\epsilon) = f(x_k) + f'(x_k)(\epsilon - x_k) + \frac{f''(\eta_k)(\epsilon - x_k)^2}{2} \quad\quad \lvert \eta_k - \epsilon \rvert \leq \lvert x_k - \epsilon \rvert
\end{equation*}
da cui
\begin{equation*}
	x_{k+1} - \epsilon = \frac{f''(\eta_k)(\epsilon - x_k)^2}{2 f'(x_k)}
\end{equation*}
e per continuità delle due derivate si ottiene
\begin{equation*}
	\lim_{k \to +\infty} \frac{\lvert x_{k+1} - \epsilon \rvert}{\lvert x_k - \epsilon \rvert^2} = \bigg\lvert \frac{f''(\epsilon)}{2f'(\epsilon)} \bigg\rvert
\end{equation*}
\end{document}
